\documentclass[11pt]{article}

\usepackage[margin=1in]{geometry}
\usepackage{amsmath,amssymb}
\usepackage{array}
\usepackage{booktabs}
\usepackage{enumitem}
\usepackage{hyperref}
\usepackage{microtype}
\usepackage[T1]{fontenc}
\usepackage{lmodern}

\hypersetup{
    colorlinks=true,
    linkcolor=blue,
    urlcolor=blue
}

\title{\textit{cs-mail}: Express Lanes for Transactional Mail\\
\large A Pedagogical Memo}
\author{C-SQD}
\date{Draft 0.3 --- September 2026}

\setlength{\emergencystretch}{1.5em}

\begin{document}
\maketitle

\begin{center}
\fbox{\parbox{0.92\textwidth}{\small\textbf{Nonnormative memo.} This document
explains a concept and motivates it for a reader new to the underlying ideas. It
introduces \emph{express lanes}, the user-facing form of the scoped capabilities
described in the \textit{cs-mail Protocol Specification}. The settlement and
privacy semantics remain governed by that specification. The companion
implementation plan develops issuance and lifetime; the deployment profile
defines the separate C-SQD financial program.}}
\end{center}

\section{A First Week with a New Bank Account}

{\sloppy
You opened a checking account on Monday. By Friday your inbox holds a welcome note
from \texttt{welcome@bank.com}, a statement notice from
\texttt{statements@bank.com}, a fraud alert from \texttt{alerts@e.bank.com}, a
marketing message you did not really ask for, and---indistinguishable at a
glance---a message from ``Bank Security Team''
\texttt{<security@\allowbreak bank-\allowbreak support-\allowbreak online.com>}
urging you to confirm your login. The
first three you wanted. The last is a phishing attempt. Every one of them arrived
the same way: someone typed your address and pressed send.\par}

\textit{cs-mail} already changed that story for strangers. An unknown person must
submit one conditionally charged relationship request with an initial message,
and you accept or reject it. Permitted follow-ups share that request's charge
and deadline.
But your bank is neither a stranger you want to charge nor a person you
``accept'' as a standing contact. You chose to hear from it, for particular
reasons, and it will write to you from many addresses for years. The two states
\textit{cs-mail} already has---requesting sender and accepted contact---do not fit
it. That gap is what this memo is about.

\section{The Problem: Two Kinds of Wanted Mail}

Wanted mail comes in two shapes. The first is people and ongoing relationships: a
colleague, a friend, a new client. \textit{cs-mail} handles these by acceptance---
you grant a person standing permission and they write freely thereafter. The
second is organizations you have authorized for a bounded purpose: the bank's
statements, the airline's boarding pass, the pharmacy's refill notice, the clinic's
appointment reminder. Call this \emph{transactional mail}.

Transactional mail does not fit acceptance, for four reasons:

\begin{itemize}[leftmargin=2em]
    \item It arrives from many addresses. One bank sends from
    \texttt{statements@}, \texttt{alerts@}, and
    \texttt{no-\allowbreak reply@\allowbreak e.bank.com}.
    Accepting an address would mean accepting dozens, and missing the next one.
    \item It arrives at volume and without your reply. You will never hand-curate
    it.
    \item It is narrow. You want the bank's fraud alerts; you did not necessarily
    invite its marketing.
    \item The sender usually is not on \textit{cs-mail} yet. Your bank still sends
    ordinary email, and will for a long time.
\end{itemize}

So we need another permission path: a way to grant a chosen organization a
bounded, revocable, effortless right to reach you---one that resists
impersonation and works \emph{before} the sender adopts \textit{cs-mail}. We call
it an \emph{express lane}. Before describing it, two pieces of background make it
comprehensible.

\section{What You Need to Know First}

\subsection{The ``From'' line is not identity}

The sender shown on an email is like the return address handwritten on an
envelope: a claim, not a proof. Anyone can write ``Your Bank'' as the sender name
and \texttt{security@bank.com} as the address; nothing in ordinary email stops
them. Phishing is precisely this move. An express lane therefore can never be
granted to a display name or a typed address, because neither is evidence of who
actually sent the message.

\subsection{Mail can prove authorized domain use: SPF, DKIM, and DMARC}

Two widely deployed standards let a receiving system check a message
cryptographically, and most newcomers have never heard of them.

\emph{DKIM} attaches a domain signature to a message; \emph{SPF} authenticates the
domain responsible for the SMTP sending path. Either can authenticate a domain,
but that domain is not necessarily the one visible to the reader.

\emph{DMARC}, standardized by RFC~9989, ties a passing SPF or DKIM identity to the
visible RFC 5322 author domain through strict or relaxed alignment. It also
lets a domain publish its requested handling for failures.

Together they answer one precise question: did the owner of the visible author
domain authorize its use on this message? They do not say the message is friendly,
wanted, legally from a particular company, or safe. A look-alike domain can pass
DMARC for itself. The lane's separate expected-domain match supplies that missing
decision.

\subsection{What \textit{cs-mail} already provides}

\textit{cs-mail} already gives us relationship requests for strangers and the idea
of recipient-granted permission. The protocol specification also already defines a
general \emph{scoped capability}---a recipient-signed grant that lets a sender
communicate without a request charge for a bounded purpose. Express lanes are that launch capability, given the
issuance, safety, and lifetime rules it needs to actually work.

\section{The Solution: Express Lanes}

An express lane is a recipient-granted, revocable, uncharged path for one sender
you have chosen, bounded in purpose and time. Three properties make it safe.

\paragraph{What it limits.}
A lane controls \emph{who} may reach you, \emph{how much}, and for \emph{how
long}. It can also constrain the authenticated purpose and automation mode that a
sender is allowed to declare. Because \textit{cs-mail} content is end-to-end
encrypted, no provider reads a message to confirm it is ``really about your
account'' or that a human wrote it. The lane enforces the sender's accountable
declaration, not the semantic truth of the ciphertext. If a sender uses that scope
dishonestly, you can report it and revoke or block the lane. This is a deliberate
limit, and an honest one.

\paragraph{What it binds to.}
For a sender already on \textit{cs-mail}, the lane binds to the sender's
cryptographic protocol identity, and the provider can honor it without even
learning which organization it is. For a legacy sender---your bank today---the
lane binds to the canonical \emph{organizational domain} established by a DMARC
pass, \texttt{bank.com}, and honors a message only when its visible author domain
falls under that identity. It never binds to a
display name or a single address. So the same lane behaves sensibly across every
case in the opening story:

\begin{itemize}[leftmargin=2em]
    \item the fraud alert from \texttt{alerts@e.bank.com} and the statement from
    \texttt{statements@bank.com} both ride the lane, because both authentically
    come from \texttt{bank.com};
    \item ``Bank Security Team''
    \texttt{<security@bank-support-online.com>} does not, because its
    authenticated domain is not \texttt{bank.com};
    \item \texttt{security@bank.com.evil.com} does not, because its authenticated
    domain is \texttt{evil.com}.
\end{itemize}

The phishing message from Friday fails at exactly this line.

\paragraph{Why it is safe to grant freely.}
An express lane is trust you can withdraw in one tap. If it is abused---the bank
floods you, or a lane is misused---you can revoke it. Future messages then need
another valid permission path. A sender with an open relationship request may
send only the follow-ups your settings allow; otherwise a new request requires
eligibility and explicit sender authorization. Revocation never silently charges
the sender or settles an existing request. A block always overrides a lane.

\medskip

These permission paths explain the main cases:

\begin{center}
\small
\begin{tabular}{@{}>{\raggedright\arraybackslash}p{0.24\textwidth}>{\raggedright\arraybackslash}p{0.30\textwidth}>{\raggedright\arraybackslash}p{0.32\textwidth}@{}}
\toprule
\textbf{Who} & \textbf{How they reach you} & \textbf{How it is granted and ended}
\\
\midrule
A stranger & Submits one charged relationship request, with any permitted
follow-ups. & Explicit sender authorization; decision or expiry settles the request. \\
A person you accept & Writes freely, no request charge. & You accept the relationship; you
revoke it. \\
An organization you chose & Travels a uncharged express lane, bounded in purpose
and time. & Issued from a decision you already made; you revoke it or let it
lapse. \\
\bottomrule
\end{tabular}
\end{center}

\paragraph{The legacy binding is a bridge.}
Binding to a domain is a bridge, not the destination. The day your bank adopts
\textit{cs-mail}, its lane re-binds to the bank's cryptographic identity, the
DMARC dependency falls away, and the lane becomes fully private---the provider no
longer needs to see \texttt{bank.com} to honor it. The bridge advertises its own
obsolescence.

\section{How You Get One: Issuance}

An express lane should feel like a consequence of a decision you already made,
never a settings screen you must remember to visit. Two moments cover almost
everything:

\begin{enumerate}[leftmargin=2em]
    \item \textbf{When you hand over your address.} Signing up, checking out, or
    booking---the moment you give the bank your address, the client offers, in
    one line, to open an express lane for it.
    \item \textbf{When you write first.} Emailing the bank yourself auto-opens a
    return lane, since you have plainly invited a reply.
\end{enumerate}

For a legacy sender the grant cannot complete on a promise; it completes on proof.
When you sign up on \texttt{bank.com}, the client remembers that you were on
\texttt{bank.com} and expects mail from that domain. The first message that
arrives claiming to be the bank is checked with DMARC; if it authentically comes
from the expected domain, the client shows a one-tap confirmation naming the
\emph{authenticated domain}---not the display name---and the lane opens. Seeding
the expectation from the website you were actually on turns ``which bank?'' into a
simple match, and lets you catch a look-alike domain at the very moment of
granting.

\section{How a Lane Ends: Lifetime and Renewal}

A lane that lasted forever by default would rebuild the very allowlist we are
trying to avoid. So lifetime is bounded, but bounded intelligently.

\paragraph{Anchor on activity, not the calendar.}
A lane's clock is its last authenticated message, not its issue date. A lane in
active use stays open on its own; one that goes quiet drifts toward expiry. This
captures most of the benefit of tying a lane to an event---a single flight,
say---without asking anyone to define events.

\paragraph{The default is a gentle re-confirm.}
When an inactive lane reaches the end of its window, you get a light prompt---keep
the airline's lane, or let it lapse?---rather than silent deletion. You set the
window per sender.

\paragraph{The exception that matters.}
Some lanes should effectively never lapse. Your bank's fraud alerts and your
hospital's test results are exactly the mail whose silent disappearance is most
costly: a lapsed lane there means a missed fraud warning or a missed result,
unable to use the lane until authority is renewed. A new charge is never automatic. For these you choose a \emph{persistent}
lane: activity renews its ordinary horizon, and an inactive horizon produces a
review reminder instead of a lapse. Even a persistent lane has a recipient-signed
hard validity bound, after which a fresh grant is required.

Revocation and blocking remain available and immediate, whatever the lifetime.

\section{Implementation Constraints}

Four constraints shape any implementation. Their details belong in a later
specification; the constraints themselves are worth stating now because they are
what make the concept real.

\paragraph{Lanes live at the recipient's provider.}
Only the provider sees inbound mail and runs the DMARC check, and the provider is
already where an incoming message is evaluated for accepted, lane, or open-request permission before any new-request flow.
The lane must be evaluated there. A device cannot verify a domain on mail it never
handled.

\paragraph{The list of lanes is sensitive and must be compartmentalized.}
Which organizations a person trusts, and for what---their bank, their hospital,
their lawyer---is precisely the relationship profile that \textit{cs-mail}'s
privacy model exists to protect. Lanes must therefore live in the relationship
privacy domain, under their own keys and access role, not in a general table that
a ledger or telemetry breach could read. For senders on \textit{cs-mail}, a lane
can be a recipient-signed, pairwise-scoped grant that the provider verifies
without learning the organization at all; for legacy senders the provider must see
the domain to match DMARC, which is one more reason the native path is the more
private destination.

\paragraph{This promotes an existing feature rather than inventing settlement.}
The protocol specification already defines scoped capabilities, already requires
that a capability-authorized message open no relationship solicitation and advance
no repeated-attempt state, and already lets a block override a capability. Express
lanes adopt those rules unchanged. The launch extension adds the issuance flow
(the two moments and the
DMARC-confirmed legacy grant), the domain-binding and expected-domain mismatch rules,
and the activity-anchored lifetime with its persistent option. The lane retains these properties under the revised request model; its grant does not create or settle a charge.

\paragraph{Declarations remain minimal.}
The authenticated envelope exposes only the coarse purpose, origin mode, validity,
and capability reference required for admission or recipient-controlled routing.
Structured message fields remain encrypted. A native sender signs its own
declarations; a legacy gateway identifies itself as the declaration authority and
does not turn DMARC into a claim about purpose or human authorship. These metadata
receive the same minimization, compartmentalization, and retention discipline as
other relationship signals.

\paragraph{Lane activity and member activity are different.}
An authenticated incoming message can renew a lane's horizon. It does not by
itself establish intentional activity for a quarterly member distribution, extend
a pending relationship request, or create a pool contribution. Those lifecycles
remain separate even when their events occur close together.

\section{Where This Leaves \textit{cs-mail}}

With express lanes, the permission model reads as one sentence: strangers submit a
single relationship request to seek permission, people you accept may write freely, and
organizations you have chosen travel an express lane you can close at any time.
Each state matches how the relationship actually feels---a stranger earning a
first hearing, a person you know, an institution you deliberately signed up with.
That is what mindful email means: not an inbox that guesses what you want, but one
where reaching you is a permission you granted, on purpose, and can always take
back.

\end{document}
